\chapter{Описание программы}
\label{chapter:program}

В данном разделе описывается функциональные возможности программы.

\section{Общие сведения}

Данная утилита была разработана в исследовательских академических целях, для
практической проверки работоспособности теоретических методов, описанных в
разделе~\ref{chapter:simulated_annealing}.

Программа является кроссплатформенной, так как реализована на языке C++,
компилятор и набор библиотек, которого имеет широкую распространённость на
различных операционных системах.

Данная реализация может в дальнейшем развиваться и агрегировать в себе другие
методы формирования и оценки S-блоков.

\section{Функциональное назначение}

Данная программа реализует работу с S-блоками, базовый математический аппарат
для исследовения характеристик S-блоков, тестирование нелинейных узлов замен
критериями, выдвигаемыми к шифру DES, а так же метод имитации отжига S-блоков
для произвольной размерности нелинейных узлов замен. Так же в программе
по-умолчанию реализованы методы оценки и формирования нелинейных узлов замен для
шифров: DES, MacGuffin, ГОСТ~28147-89.

\section{Описание логической структуры}

На рисунке~\ref{fig:program_logic} продемонстрирован алгоритм работы программы в
режиме имитационного моделирования. В каждом цикле тестирования случайно
инициализируется S-блок, после чего применяется алгоритм имитации отжига,
описанный в разделе~\ref{chapter:simulated_annealing}, и в результате
оцениваются характеристики полученного S-блока.

Основными функциями программы являются:
\begin{itemize}

    \item Синтез нелинейных узлов замен методом имитации отжига
    (см. приложение~\ref{app:listing_simulated_annealing}): \\
    void SimulatedAnnealing::run()

    \item Оценка нелинейности, автокорреляции S-блока (см. приложение~\ref{app:listing_NL_AC}): \\ 
    int SBox::get\_NL(), \\
    int SBox::get\_AC();

    \item Оценка верхних границ вероятности дифференциального и линейного
    криптоанализа ГОСТ~28147-89 (см. приложение~\ref{app:listing_MD_ML}): \\
    double SBox::get\_MD(), \\
    double SBox::get\_ML();

\end{itemize}

\begin{figure}
    \centering\includegraphics[width=90mm]{program_logic.eps}
    \caption{Алгоритм работы программы}
    \label{fig:program_logic}
\end{figure}


\section{Используемые технические средства}

Для работы данной программы необходимо устройство с процессором, архитектура
которого поддерживается компилятором С++ кода (например, GCC поддерживает все
наиболее распространённые архитектуры и множество менееизвестных). Реализация
протестирована и гарантированно функционирует в ОС GNU/Linux. 

\section{Вызов и загрузка}

Программа не принимает никаких параметров в момент запуска, все параметры
константно задаются перед компиляцией.

\section{Входные и выходные данные}

Входными данными является случайно сгенерированное начальное состояние
нелинейного узла замен.

Выходными данными является синтезированный нелинейный узел замен и его
характеристики. Так же программа имеет режим тестирования, в котором она
вычисляет наихудшие и средние характеристики сформированных S-блоков по
критериям нелинейности, автокорреляции и значениям таблицы разностей.
